% Pursuing Stacks (cotes 134-2 to 134-8): archive cote 134-6, site page 46 (PDF page 47), typescript page 364.
% typescript. Draft compared: ps5.tex:3633-3711.
% AI-transcribed by Claude Opus 5.5 (claude-opus-5-5, adaptive thinking, effort high), Message Batches API, batch ps3 (msgbatch_01XdFmuVVBAcgc7EFxdtCqsj), 2026-11-12; reviewed against both supplied scan halves; see CHECK_DECISIONS.tsv.
\scanpage{134-6/46}
\mspage{364}

We'll henceforth identify Add(B) to a full subcategory of $B^{\hand{\wedge}}_{ab}$
(by replacing the solution of the universal problem, constructed in
section 93, by the essential \xout{dirx} image in $B^{\hand{\wedge}}_{ab}$ say), and rewrite (5)
simply as
\[ (5') \qquad \overset{\hand{\sim}}{F}(L)\hand{\simeq} \mathrm{Hom}(\ L\ ,\ F) \quad , \]
the Hom being taken in $B^{\hand{\wedge}}_{ab}$ , category of abelian presheaves on B.
Accordingly, if L. is a \over{c}{c}\over{h}{h}ain complex in Add(B), hence in $B^{\hand{\wedge}}_{ab}$ ,
the corresponding \uline{cochain} complex $\overset{\hand{\sim}}{F}(L.)$ in (Ab) can be interpreted as
\[ (6) \qquad \overset{\hand{\sim}}{F}(L.)\hand{\simeq} \mathrm{Hom}^{\bullet}(L.,F) \quad , \]
where \xout{the} the symbol Hom$^{\bullet}$ means taking Hom's componentwise.

\tikz[baseline=(ww.base)]{\node[inner sep=0pt] (ww) {\xout{Wh\unsure{w}at}};\draw ([xshift=-1em,yshift=.7ex]ww.north) .. controls ([xshift=-.65em,yshift=-.1ex]ww.north) .. ([xshift=-.3em,yshift=.7ex]ww.north);\draw ([xshift=-.95em,yshift=-.65ex]ww.south) .. controls ([xshift=-.65em,yshift=-.1ex]ww.south) .. ([xshift=-.35em,yshift=-.65ex]ww.south);} What we're after here is to find a \uline{fixed} chain complex L. in
Add(B), such that for \uline{any} abelian presheaf F on B, the cochain com-%
plex (6) in (Ab) should be \over[typed]{a}{i}somorphic (in the derived category \xout{D(Ab)}
$D^{\bullet}$(Ab) of cochain complexes in (Ab) with respect to quasi-isomorphisms)
to the ``integration'' of F over the topos $B^{\hand{\wedge}}$ , i.e. to $R\hand{\Gamma}_B(F)$ :
\[ (*) \qquad \mathrm{Hom}^{\bullet}(L.,F)\hand{\simeq} R\hand{\Gamma}_B(F) \quad ? \quad (\text{isom. in } D^{\bullet}(\mathrm{Ab}))\ , \]
namely to the total right derived functor $R\hand{\Gamma}_B$ (taken \xout{on F} for the argu%
ment F) of the ``sections'' functor
% NOTE: the end of the line "argu" is cut at the right edge of the scan (a hyphen may be lost)
\margin{\tikz[baseline=(lm.base)]{\node[inner sep=4pt] (lm) {\ensuremath{\displaystyle\varprojlim_{B^{\circ}} F}};\draw[color=blue] (lm.center) ellipse (2.7em and 2em);}}
\[ (7) \qquad \hand{\Gamma}_B(F) \overset{\hand{def}}{=} \underset{\over{B\unsure{B}}{B^{\circ}}}{\underset{\hand{\longleftarrow}}{\mathrm{lim}}}\ F \quad . \]
Now, using the fact that the components of the chain complex L. are
projective, hence $\mathrm{Ext}^{i}(L_n,F) = 0$ for $i > 0$ (any n, any F), we get
at any rate a canonical isomorphism in $D^{\bullet}$(Ab), or \inserted[typed]{in} D(Ab) :
\[ (8) \qquad \mathrm{Hom}^{\bullet}(L.,F) \hand{\simeq} R\tikz[baseline=(rh.base)]{\node[inner sep=0pt] (rh) {\ensuremath{\unsure{h}}};\node[inner sep=0pt] at (rh.center) {\ensuremath{H}};}\mathrm{om}(L.,F) \quad , \]
i.e. an interpretation of (6) as an ``hyperext'' . Now let's remember
that $R\hand{\Gamma}_B(F)$ (as on any topos) can be interpreted equally as
\[ (9) \qquad R\hand{\Gamma}_B(F)\hand{\simeq} R\mathrm{Hom}(\underline{Z}_B,F) \quad , \]
where $\underline{Z}_B$ denotes the constant presheaf on B with value $\underline{Z}$ . Thus \over{a}{t}he
wished-for isomorphism (*) will follow most readily from a correspon-%
ding isomorphism in $D.(B^{\hand{\wedge}}_{ab})$ between L. and $\underline{Z}_B$ . But using again
the fact that the components of L. are projective, we see that to give
a map in the derived category of L. into $\underline{Z}_B$ amounts to the same as to
give an \uline{augmentation}
\[ (10) \qquad L. \hand{\longrightarrow} \underline{Z}_B \quad , \]
and the map is an isomorphism in $D.(B^{\hand{\wedge}}_{ab})$

\margin{46}

% Visual collation: both supplied image halves inspected; see CHECK_DECISIONS_134_7_8.tsv.
